% ============================================================
%  Restauro de Passé — Manual de regras
% ============================================================
\documentclass[11pt]{article}
\usepackage[a4paper,margin=18mm]{geometry}
\usepackage[T1]{fontenc}
\usepackage[utf8]{inputenc}
\usepackage[portuguese]{babel}
\usepackage{lmodern}
\usepackage{microtype}
\usepackage{xcolor}
\usepackage{tikz}
\usetikzlibrary{calc}
\usepackage{multicol}

% ---------- Paleta (idêntica às cartas) ----------
\definecolor{fachada}{HTML}{B04A2E}
\definecolor{nave}{HTML}{2E6E8C}
\definecolor{capela}{HTML}{6C4B99}
\definecolor{taipa}{HTML}{7A5A33}
\definecolor{telhado}{HTML}{3F7E58}
\definecolor{historia}{HTML}{2E5E8C}
\definecolor{construcao}{HTML}{4E7C31}
\definecolor{acao}{HTML}{B03A2E}
\definecolor{desafio}{HTML}{C4841D}
\definecolor{ouro}{HTML}{E9B44C}
\definecolor{tinta}{HTML}{333333}
\definecolor{areia}{HTML}{F2EDE1}
\definecolor{calcario}{HTML}{FAF6EC}

\usepackage{helvet}
\setlength{\parindent}{0pt}
\setlength{\parskip}{4pt}
\pagestyle{empty}

% ---------- Título de seção desenhado em TikZ ----------
\newcommand{\regrasec}[2]{%
  \par\vspace{9pt}
  \begin{tikzpicture}
    \fill[#1,rounded corners=1.5mm] (0,0) rectangle (174mm,7.5mm);
    \node[anchor=west,font=\sffamily\bfseries\large,text=white] at (3mm,3.75mm) {#2};
  \end{tikzpicture}\par\vspace{3pt}}

\newcommand{\regrasub}[1]{\par\vspace{5pt}{\sffamily\bfseries\large\color{tinta}#1}\par\vspace{1pt}}

% ---------- Caixa de destaque ----------
\newcommand{\caixa}[2]{%
  \par\vspace{4pt}
  \begin{tikzpicture}
    \node[fill=#1,rounded corners=2mm,inner sep=3.5mm,text width=164mm] {\footnotesize #2};
  \end{tikzpicture}\par\vspace{2pt}}

\newcommand{\dica}[1]{\caixa{areia}{{\sffamily\bfseries\color{tinta}Dica de mestre de obras:} #1}}

% ---------- Peças da igreja ----------
\newcommand{\partebotao}[3]{% cor, inicial, nome
  \begin{tikzpicture}[baseline=-0.5ex]
    \node[circle,fill=#1,minimum size=5mm,inner sep=0pt,font=\sffamily\bfseries\footnotesize,text=white] {#2};
    \node[anchor=west,font=\footnotesize\sffamily] at (3.5mm,0) {#3};
  \end{tikzpicture}}

\begin{document}

% ===================== CAPA =====================
\begin{tikzpicture}[remember picture,overlay]
  \fill[calcario] (current page.south west) rectangle (current page.north east);
  \fill[areia] ([yshift=200mm]current page.south west) rectangle (current page.north east);
  % fachada
  \draw[color=historia,line width=2pt,line join=round]
    ([xshift=45mm,yshift=150mm]current page.south west) -- ([xshift=75mm,yshift=150mm]current page.south west)
    -- ([xshift=75mm,yshift=205mm]current page.south west)
    -- ([xshift=105mm,yshift=225mm]current page.south west)
    -- ([xshift=135mm,yshift=205mm]current page.south west)
    -- ([xshift=135mm,yshift=150mm]current page.south west);
  % porta em arco
  \draw[color=historia,line width=2pt]
    ([xshift=85mm,yshift=150mm]current page.south west) -- ([xshift=85mm,yshift=173mm]current page.south west)
    arc[start angle=180,end angle=0,radius=10mm]
    -- ([xshift=105mm,yshift=150mm]current page.south west);
  % cruz
  \draw[color=historia,line width=2pt]
    ([xshift=105mm,yshift=225mm]current page.south west) -- ([xshift=105mm,yshift=237mm]current page.south west)
    ([xshift=100mm,yshift=231mm]current page.south west) -- ([xshift=110mm,yshift=231mm]current page.south west);
  % torre sineira
  \draw[color=historia,line width=2pt,line join=round]
    ([xshift=33mm,yshift=150mm]current page.south west) -- ([xshift=33mm,yshift=208mm]current page.south west)
    -- ([xshift=47mm,yshift=208mm]current page.south west) -- ([xshift=47mm,yshift=150mm]current page.south west)
    ([xshift=30mm,yshift=208mm]current page.south west) -- ([xshift=40mm,yshift=217mm]current page.south west)
    -- ([xshift=50mm,yshift=208mm]current page.south west);
\end{tikzpicture}

\vspace*{168mm}
\begin{center}
  {\sffamily\bfseries\Huge\color{tinta} RESTAURO DE PASSÉ\par}
  \vspace{4mm}
  {\sffamily\Large\color{tinta} O jogo de cartas da Igreja de Nossa Senhora\\da Encarnação de Passé --- Candeias, Bahia\par}
  \vspace{7mm}
  {\sffamily\large\color{historia} Manual de regras para guardiões de 6 a 14 anos\par}
  \vspace{12mm}
  {\sffamily\small\color{tinta} Um jogo educativo sobre patrimônio histórico e\\técnicas construtivas do Recôncavo Baiano\par}
\end{center}
\newpage

% ===================== 1. O QUE É =====================
\regrasec{historia}{O que é este jogo?}

\textbf{Restauro de Passé} é um jogo de cartas educativo sobre a ruína da
\textbf{Igreja de Nossa Senhora da Encarnação de Passé}, no distrito de Passé,
em Candeias (Bahia). Ao jogar, as crianças e jovens conhecem a história do
monumento e as \textbf{técnicas construtivas} usadas pelos antigos mestres de
obra --- taipa de pilão, pedra e cal, telhas coloniais --- que ainda podem
inspirar as construções atuais do povoado.

Cada carta tem um \textbf{SABER}: um fato real sobre a igreja ou sobre uma
técnica de construção. Ler em voz alta faz todo mundo aprender, mesmo quando
não é a sua vez.

\caixa{areia}{\textbf{Para quem é:} crianças e jovens de 6 a 14 anos, em casa,
na escola ou em oficinas comunitárias. Pode ser jogado por leigos e usado
como material didático em cursos de arquitetura e patrimônio.}

% ===================== 2. CONTEÚDO =====================
\regrasec{tinta}{O que vem no jogo?}

O baralho tem \textbf{54 cartas} em quatro naipes:

\begin{itemize}
  \item \textcolor{historia}{\textbf{16 cartas HISTÓRIA}} (azul): fatos sobre a
  igreja, o povoado de Passé e a vida de quem ali viveu.
  \item \textcolor{construcao}{\textbf{16 cartas CONSTRUÇÃO}} (verde): as
  técnicas e os materiais usados para erguer a igreja --- e como reaproveitá-los hoje.
  \item \textcolor{desafio}{\textbf{12 cartas DESAFIO}} (laranja): perguntas
  sobre o monumento. Acertou, ficou com a carta!
  \item \textcolor{acao}{\textbf{10 cartas AÇÃO}} (vermelho): acontecimentos e
  brincadeiras que mudam o rumo do jogo.
\end{itemize}

Cada carta de HISTÓRIA e de CONSTRUÇÃO pertence a uma \textbf{peça da
igreja}, marcada no canto inferior direito:

\begin{center}
\partebotao{fachada}{F}{Fachada e torre}\quad
\partebotao{nave}{N}{Nave}\quad
\partebotao{capela}{C}{Capela-mor}\quad
\partebotao{taipa}{T}{Paredes}\quad
\partebotao{telhado}{S}{Telhado}\quad
\partebotao{ouro}{$\star$}{Curinga (DESAFIO)}
\end{center}

(dica: ``S'' vem de \emph{sobrado} --- ninguém quer que a chuva molhe o forro!)

% ===================== 3. OBJETIVO =====================
\regrasec{telhado}{O objetivo: reconstruir a igreja!}

Você vence quando reúne, na sua frente, \textbf{cartas das cinco peças da
igreja}: Fachada, Nave, Capela-mor, Paredes e Telhado.

As cartas DESAFIO (curinga) podem ser usadas como \textbf{qualquer peça} que
estiver faltando. Você decide na hora!

\dica{As cartas valem pontos (o número amarelo no alto). Se duas pessoas
reconstruírem a igreja na mesma rodada, vence quem tiver mais pontos.}

% ===================== 4. PREPARO =====================
\regrasec{fachada}{Preparando o jogo}

\begin{enumerate}
  \item \textbf{Separe} o baralho em quatro montinhos pelos naipes (cores).
  \item \textbf{Embaralhe} todos os montinhos juntos e forme o monte de
  compra, no centro, com o verso para cima.
  \item Cada jogador \textbf{compra 3 cartas}. Não mostre para ninguém!
  \item Os jogadores se organizam em círculo. O mais novo começa.
\end{enumerate}

% ===================== 5. COMO JOGAR =====================
\regrasec{capela}{Como se joga}

O jogo anda no sentido horário. Na sua vez:

\begin{enumerate}
  \item \textbf{Compre} 1 carta do monte.
  \item \textbf{Leia em voz alta} o texto SABER da carta, para todo mundo
  aprender junto.
  \item \textbf{Cumpra a carta}, conforme o naipe:
\end{enumerate}

\regrasub{Cartas HISTÓRIA e CONSTRUÇÃO (azuis e verdes)}
Guarde a carta na sua frente, junto das outras da mesma peça da igreja. Se
você tiver as \textbf{cinco peças}, grite ``\textbf{A igreja está de pé!}'' ---
você venceu!

\regrasub{Cartas DESAFIO (laranja)}
O jogador à sua \textbf{direita} lê a pergunta em voz alta (sem ler a
resposta, claro!).
\begin{itemize}
  \item \textbf{Acertou:} fique com a carta. Ela é um curinga: use como a peça que quiser.
  \item \textbf{Errou:} a carta vai para o fundo do monte de compra.
\end{itemize}

\regrasub{Cartas AÇÃO (vermelhas)}
Faça o que a carta manda \emph{na hora} --- e leia o SABER delas também, que
cada uma conta uma história.

\regrasub{Fim da vez}
Se o monte de compra acabar, reembaralhe as cartas descartadas (menos as que
estão na frente dos jogadores) e continue jogando.

% ===================== 6. AÇÕES =====================
\regrasec{acao}{As cartas de AÇÃO}

Cada carta vermelha explica o próprio efeito no texto. Em geral, elas fazem
você:

\begin{itemize}
  \item \textbf{Pilão!} --- comprar cartas extras e distribuir SABER;
  \item \textbf{Troca-troca} --- trocar cartas com outro jogador;
  \item \textbf{Mestre de obras} --- ``roubar'' uma carta de peça que está
  faltando;
  \item \textbf{Chuva forte} --- perder a vez e proteger as paredes de taipa;
  \item \textbf{Romaria} --- todos os jogadores ganham uma carta e leem um
  SABER juntos.
\end{itemize}

Se uma carta mandar pegar cartas de outro jogador e ele não tiver, nada
acontece --- sem brigas de terreiro!

% ===================== 7. IDADES =====================
\regrasec{historia}{Adaptando para cada idade}

\regrasub{6 a 8 anos --- pequenos ajudantes de obra}
\begin{itemize}
  \item Jogue \textbf{em dupla}, criança + adulto ou dois colegas, dividindo as cartas.
  \item Nas perguntas, leia \textbf{duas alternativas em voz alta} e deixe a
  criança escolher.
  \item Valide qualquer tentativa de ler o SABER: errar a palavra não importa,
  quem importa é a vontade de contar a história.
\end{itemize}

\regrasub{9 a 11 anos --- aprendizes de mestre de obras}
\begin{itemize}
  \item Jogo normal, com todas as regras.
  \item Crie um \textbf{caderno de obra}: anote um fato novo por partida.
\end{itemize}

\regrasub{12 a 14 anos --- mestres de obras}
\begin{itemize}
  \item Modo \textbf{pontuação}: além das cinco peças, some os pontos das
  cartas. Precisa ter as cinco peças \textbf{e} pelo menos 20 pontos para vencer.
  \item Ao final, cada vencedor conta \textbf{de memória} um fato da igreja
  para a turma --- contou errado, perde 2 pontos!
\end{itemize}

% ===================== 8. COOPERATIVO =====================
\regrasec{construcao}{Modo cooperativo --- todos contra a chuva}

Para turmas e famílias que preferem construir juntos:

\begin{enumerate}
  \item Não há vencedor individual: o grupo vence \textbf{junto}.
  \item As cartas ficam no centro, organizadas por peça da igreja.
  \item O grupo precisa montar as \textbf{cinco peças antes do monte acabar
  duas vezes} (na segunda passada, reembaralhe só uma vez).
  \item Conseguiram? A igreja está de pé! Não conseguiram? A chuva venceu ---
  recomecem, agora sabendo mais do que antes.
\end{enumerate}

% ===================== 9. GLOSSÁRIO =====================
\regrasec{taipa}{Glossário das técnicas (para aprendizes e arquitetos)}

{\footnotesize
\begin{itemize}
  \item \textbf{Taipa de pilão} --- parede feita de terra úmida socada dentro de
  um molde de madeira (a taipa), levantada em camadas de uns 40\,cm. Quando a
  camada seca, o molde sobe. É forte, térmica e usa material do próprio lugar.
  \item \textbf{Adobe (tijolo crú)} --- tijolos de barro moldados ao sol, em vez de
  serem queimados. Mais leve que a taipa e fácil de remendar.
  \item \textbf{Pedra e cal} --- argamassa feita de cal (pedra calcária
  queimada) misturada a areia e pedrisco. Cola as pedras e recebe o reboco.
  \item \textbf{Reboco} --- camada fina de cal e areia que protege a parede da
  chuva e deixa as superfícies lisas para pintura.
  \item \textbf{Telha colonial} --- telhas de barro canal e capa, moldadas na
  coxa ou em fôrma, que fazem o telhado de água correr para fora e o forro
  respirar.
  \item \textbf{Madeiramento} --- estrutura de madeira que segura o telhado:
  tesouras, caibros e ripas, quase sempre de madeira local.
  \item \textbf{Verga e contra-verga} --- peças de madeira ou pedra que
  seguram o peso da parede acima de portas e janelas, sem deixar rachar.
  \item \textbf{Capela-mor} --- a parte da igreja onde fica o altar, geralmente
  mais cuidada e às vezes com pé-direito maior que a nave.
  \item \textbf{Nave} --- o corpo principal da igreja, onde ficam as pessoas
  durante as celebrações.
\end{itemize}
}

% ===================== 10. OFICINAS =====================
\regrasec{telhado}{Do jogo para a obra: oficinas práticas}

O baralho é porta de entrada. Depois de jogar, leve as técnicas para a vida
real --- é assim que a igreja ensina as gerações de hoje:

\regrasub{Oficina 1 --- Minha primeira taipa}
Material: caixa de sapato aberta nas laterais, terra úmida e areia, um pedaço
de pau como pilão. Encha a caixa de 4 em 4 dedos e soca bem, como na carta da
taipa. Ao final, compare com uma parede de taipa verdadeira (fotos ajudam).

\regrasub{Oficina 2 --- Caça-técnica no bairro}
Em grupos, procurem no bairro casas com taipa, adobe, reboco de cal ou telha
colonial. Foto + carta correspondente = ponto do grupo. \textbf{Regra de
ouro:} olhar e fotografar só com permissão do morador; nunca tocar ou danificar nada.

\regrasub{Oficina 3 --- Minha casa, minha técnica}
Desenhe a fachada da sua casa (ou da casa dos seus avós) trocando os
materiais por técnicas da igreja: e se sua casa tivesse paredes de taipa?
E o telhado colonial? Apresente o desenho para a turma contando por que a
técnica funciona no clima daqui.

% ===================== 11. CUIDADO =====================
\regrasec{acao}{Cuidado! Isto é patrimônio}

A Igreja de Passé é um bem de todos. A ruína \textbf{não deve ser escalada},
não se deve retirar pedras, telhas ou argila do local, e visitas devem ser
feitas com respeito --- de preferência acompanhadas de moradores e guias
locais. O jogo é convite para \textbf{conhecer e proteger}, nunca para
perturbar o monumento.

% ===================== 12. CRÉDITOS =====================
\regrasec{tinta}{Créditos e licença}

\textbf{Restauro de Passé} --- jogo de cartas educativo em LaTeX/TikZ sobre a
Igreja de Nossa Senhora da Encarnação de Passé (Candeias, Bahia).

\begin{itemize}
  \item Projeto aberto: texto e arte das cartas sob licença
  \textbf{CC BY-SA 4.0}; código-fonte LaTeX sob \textbf{MIT}.
  \item Fontes históricas consultadas estão em \texttt{docs/FONTES.md}.
  \item Impressão: papel A4, 100\,\% de escala (sem ``ajustar à página'').
  Cartas de 63\,$\times$\,88\,mm.
  \item Contribuições (novas cartas, correções, traduções) são bem-vindas ---
  veja \texttt{CONTRIBUTING.md}.
\end{itemize}

\end{document}
